\chapter*{How to Finish the Syllabus in 24 Hours}
\markboth{How to Finish in 24 Hours}{}
\addcontentsline{toc}{chapter}{How to Finish the Syllabus in 24 Hours}

These notes contain \textbf{only what you write in the exam}: definitions, key points, formulas, algorithm steps, comparison tables, diagrams to draw and the standard numericals. No code and no long stories: for those, use the main textbook.

\section*{The 24-hour plan}
\begin{center}
\small
\begin{tabular}{@{}clcp{7.4cm}@{}}
\toprule
\textbf{Session} & \textbf{Chapter} & \textbf{Hours} & \textbf{Must be able to do}\\
\midrule
1 & 1 Introduction to ML & 2.0 & define ML, types of ML, challenges, overfitting\\
2 & 2 Statistical learning & 2.5 & loss, risk, bias--variance, CV, metrics, standard error\\
3 & 3 k-NN, trees, Naive Bayes & 2.5 & k-NN by hand, entropy and gain, Bayes by hand\\
4 & 4 Linear models, GLM, SVM & 2.5 & least-squares line, $R^2$, sigmoid, log-loss, margin\\
5 & 5 Ensemble learning & 2.5 & voting, bagging, random forest, AdaBoost, stacking\\
6 & 6 SVM and Naive Bayes & 2.0 & kernels, $C$ and $\gamma$, SVR, NB variants\\
7 & 7 Clustering & 2.5 & K-Means by hand, DBSCAN, GMM and EM\\
8 & 8 Dimensionality reduction & 2.0 & curse of dimensionality, PCA by hand\\
9 & 9 Neural networks with Keras & 2.5 & perceptron, MLP, activations, backpropagation\\
10 & 10 TensorFlow and data & 1.5 & tensors, \texttt{tf.data}, preprocessing layers\\
11 & Last-hour sheet + mock test & 1.5 & formulas, comparisons, one full question paper\\
\midrule
 & \textbf{Total} & \textbf{24.0} & \\
\bottomrule
\end{tabular}
\end{center}

\textbf{Suggested split:} three days of 8 hours (Day~1: Ch.~1--3 + half of 4; Day~2: rest of 4 to 7; Day~3: 8--10 and revision), or six days of 4 hours.

\section*{How to study each chapter (in this order)}
\begin{enumerate}
  \item \textbf{Definitions} (blue boxes): read aloud, then write each one from memory once. They are the 2-mark answers.
  \item \textbf{Key points and tables}: these are the body of 5-mark answers.
  \item \textbf{Diagrams to draw} (green boxes): draw each one once on paper.
  \item \textbf{Numerical patterns} (orange boxes): cover the solution and solve it yourself.
  \item \textbf{Important questions}: answer the 2-mark ones orally; outline two of the 10-mark ones.
\end{enumerate}

\section*{Answer-writing formula}
\begin{center}
\small
\begin{tabular}{@{}lp{12.0cm}@{}}
\toprule
\textbf{2 marks} & definition (2--3 lines) $+$ one formula or example\\
\textbf{5 marks} & definition $+$ diagram $+$ 4--5 points $+$ small example\\
\textbf{10 marks} & intro $+$ definition $+$ steps/algorithm $+$ diagram $+$ numerical $+$ advantages/disadvantages $+$ applications $+$ conclusion\\
\textbf{Numerical} & Given $\to$ formula $\to$ substitution $\to$ step-by-step working $\to$ answer with its meaning\\
\bottomrule
\end{tabular}
\end{center}

\begin{imp}
The boxes mean: \textbf{Definition} = write exactly; \textbf{Exam focus} = most-asked point; \textbf{Diagram to draw} = practise once; \textbf{Numerical pattern} = solve once without looking. The line \textbf{Advantages / Disadvantages / Applications} under each method is ready to copy at the end of a long answer.
\end{imp}
